\section*{ЗАКЛЮЧЕНИЕ}
\phantomsection
\addcontentsline{toc}{section}{Заключение}

В результате выполнения курсовой работы решена задача численного анализа влияния дискретного микрорельефа на трибологические характеристики радиального подшипника скольжения. Были выполнены следующие задачи:

\begin{enumerate}
    \item Проведён обзор теоретических основ гидродинамической смазки и современных исследований в области лазерного текстурирования поверхностей трения. Установлено, что ключевым механизмом положительного влияния микрорельефа является микрогидродинамический эффект, создаваемый локальными клиновыми зазорами на кромках углублений.

    \item Освоена методика численного решения уравнения Рейнольдса для подшипника конечной длины методом конечных разностей с итерационной процедурой Гаусса--Зейделя и верхней релаксацией. Расчёт выполнен на сетке $500 \times 500$ узлов с критерием сходимости $\delta < 10^{-5}$.

    \item Получены поля давления и интегральные характеристики (нагрузочная способность $F$, коэффициент трения $\mu$, расход смазки $Q$) для гладкого подшипника и подшипника с плоскодонными углублениями с фаской (тип~8) в диапазоне эксцентриситетов $\varepsilon = 0{,}05 - 0{,}80$. Геометрия подшипника: $R = 35$~мм, $c = 50$~мкм, $L = 56$~мм.

    \item Выполнен сравнительный анализ гладкого и текстурированного подшипников. При характерном рабочем эксцентриситете $\varepsilon = 0{,}6$ зафиксированы следующие изменения:
    \begin{itemize}
        \item нагрузочная способность увеличилась на $19{,}09$\%: с $F = 5234{,}7$~Н до $F = 6234{,}0$~Н;
        \item коэффициент трения снизился на $17{,}69$\%: с $\mu = 0{,}01480$ до $\mu = 0{,}01218$;
        \item расход смазки возрос на $1{,}36$\%: с $Q = 0{,}2356$~л/с до $Q = 0{,}2388$~л/с.
    \end{itemize}
\end{enumerate}

Полученные результаты подтверждают, что нанесение плоскодонных углублений с фаской с параметрами $a = 2{,}42$~мм, $b = 2{,}22$~мм, $h_p = 10$~мкм ($h_p/c = 0{,}2$) приводит к улучшению рабочих характеристик подшипника. Одновременное повышение нагрузочной способности и снижение коэффициента трения при незначительном росте расхода смазки свидетельствует об эффективности данного типа микрорельефа.

Плоскодонный профиль с фаской характеризуется постоянной глубиной в центральной зоне углубления и плавным переходом к поверхности на краях. Такая геометрия обеспечивает формирование устойчивого дополнительного давления за счёт микрогидродинамического эффекта на входной кромке фаски. Постоянная глубина центральной зоны способствует равномерному распределению смазки внутри углубления, а плавный переход на периферии исключает резкие скачки давления и разрывы потока.

На основании проведённого анализа можно рекомендовать применение плоскодонного микрорельефа с фаской для подшипников скольжения нефтегазового оборудования, работающего в условиях жидкостного трения. Особенно перспективным является использование текстурирования в тяжелонагруженных узлах ($\varepsilon > 0{,}5$), где прирост нагрузочной способности и снижение трения выражены наиболее значительно.
